\chapter{Introduzione}
\label{chap:introduction}

Il presente documento descrive le attività svolte durante il periodo di tirocinio aziendale presso FlashStart, sotto la supervisione del Prof. Mirko Viroli, i tutor aziendali Francesco Collini e Stefano Babini.

FlashStart è una azienda che da oltre 20 anni protegge gli utenti nel mondo digitale con soluzioni di sicurezza informatica attraverso il filtraggio dei contenuti basato su \ac{DNS} e intelligenza artificiale.
L'azienda fornisce filtraggio dei contenuti per utenti singoli e organizzazioni di ogni dimensione.
FlashStart opera su una rete anycast\footnote{Anycast è una metodologia di indirizzamento di rete dove le richieste vengono indirizzate ad una varietà di diversi nodi. Tipicamente si indirizza il traffico in entrata al server più vicino che ha la capacità di processare la richiesta efficientemente.} globale composta da più di 160 nodi distribuiti strategicamente in tutti i continenti, progettata per garantire elevata disponibilità, latenza minima e continuità operativa anche in caso di interruzioni o picchi di traffico.

A partire dal 2026 FlashStart ha iniziato il più grande redesign dell'azienda, spostandosi da una applicazione con architettura monolitica ad un applicativo moderno con una architettura basata su microservizi altamente scalabili.
In questo documento si farà riferimento ai due applicativi come: \textit{FlashStart Cloud}, il sistema attualmente in uso ma destinato alla dismissione, e \textit{FlashStart Internet Protection 2026}, la soluzione individuata per la sua sostituzione.

Nel seguito della trattazione, per semplicità, si farà riferimento al nuovo sistema come \textit{FlashStart 2026}.

\paragraph{Contributo della tesi}
Lo strumento di controllo della latenza rappresenta un strumento fondamentale nel toolkit del Network Admin, poiché consente di stimare a priori il tempo di risposta del servizio \ac{DNS} (espresso in millisecondi) che l'utente finale sperimenterà nell'utilizzo del servizio di protezione della navigazione offerto da FlashStart. Questa stima preventiva permette di individuare eventuali criticità infrastrutturali prima che queste si traducano in un degrado percepibile dell'esperienza utente, fornendo così un supporto decisionale concreto per le attività di monitoraggio e ottimizzazione della rete.
Dal momento che la protezione erogata da FlashStart opera ``at \ac{DNS} levek'', una latenza contenuta, tipicamente inferiore ai $10 \text{ ms}$, risulta essenziale per coniugare due esigenze apparentemente in tensione tra loro: da un lato garantire un livello adeguato di sicurezza nella navigazione, dall'altro preservarne la fluidità, evitando che l'utente percepisca rallentamenti o ritardi nella risoluzione dei nomi di dominio. Un servizio di sicurezza percepito come ``lento'' rischia infatti di essere disabilitato o aggirato dagli utenti stessi, vanificando l'obiettivo protettivo per cui è stato progettato.
Oltre alla misurazione della latenza sul nodo preferenziale, lo strumento è in grado di stimare eventuali nodi alternativi disponibili in caso di malfunzionamento o interruzione di servizio (outage) del nodo primario, individuato secondo il criterio del best latency path. Questa capacità predittiva è resa possibile dall'architettura di rete sottostante, basata sul protocollo \ac{BGP} in configurazione Anycast: gli indirizzi \ac{IP} del servizio di \ac{DNS} sicuro vengono infatti annunciati contemporaneamente da ogni datacenter della rete distribuita. Ne consegue che, in presenza di una criticità sul nodo primario, il traffico dell'utente viene automaticamente e trasparentemente rediretto verso un nodo alternativo, garantendo la continuità della navigazione filtrata e protetta, seppur con un incremento marginale della latenza rispetto alla condizione ottimale.

Il contributo principale di questo lavoro di tesi consiste nella migrazione del servizio di misurazione della latenza dalla precedente implementazione, originariamente sviluppata utilizzando i linguaggi PhP e Python, verso una nuova piattaforma progettata secondo criteri architetturali più moderni, coerenti e strutturati. Tale processo di migrazione non si è tuttavia limitato a una semplice traduzione o riscrittura del codice preesistente in un nuovo linguaggio o framework: ha comportato, al contrario, una revisione approfondita e critica dell'architettura sottostante, con l'obiettivo esplicito di superare i limiti strutturali e funzionali riscontrati nella soluzione originaria.
In particolare, il lavoro dovrà conseguire un miglioramento significativo lungo tre direttrici fondamentali. La prima riguarda l'affidabilità del sistema, intesa come la capacità di garantire un funzionamento stabile, coerente e resiliente anche in condizioni di carico elevato o di traffico anomalo, minimizzando il rischio di malfunzionamenti o interruzioni del servizio. La seconda riguarda l'efficienza operativa, perseguita attraverso l'ottimizzazione dell'utilizzo delle risorse computazionali disponibili e la conseguente riduzione dei tempi di risposta del sistema, elemento particolarmente critico trattandosi di uno strumento di misurazione della latenza stessa. La terza direttrice, infine, concerne l'usabilità della piattaforma, considerata sia dal punto di vista degli utenti finali che interagiscono con lo strumento, sia da quello degli sviluppatori incaricati di mantenerlo, estenderlo e aggiornarlo nel tempo, garantendo così la sostenibilità del progetto anche in una prospettiva di lungo periodo.
\paragraph{Struttura dei Capitoli}

\begin{description}
    \item[\Cref{chap:background} -- Sfondo tecnologico e architettura:] introduce il contesto tecnologico di riferimento, fornendo una descrizione dettagliata dell'architettura del servizio e delle sue componenti fondamentali (\emph{upstream orchestrator}, \emph{downstream} e il punto d'ingresso su FlashStart Cloud). Una particolare attenzione è rivolta all'individuazione e all'analisi sistematica delle fragilità e delle criticità dell'infrastruttura esistente, che costituiscono il punto di partenza motivazionale per il lavoro svolto.
    
    \item[\Cref{chap:analysis} -- Requisiti ed obiettivi:] definisce in modo formale i requisiti e gli obiettivi del progetto, strutturando i vincoli funzionali e architetturali prescritti per ciascun elemento del sistema sulla base delle criticità emerse dall'analisi precedente.
    
    \item[\Cref{chap:design} -- Scelte di design:] illustra le decisioni architetturali e di progettazione adottate al fine di soddisfare i requisiti individuati, formalizzando e motivando le scelte strutturali effettuate.
    
    \item[\Cref{chap:development} -- Sviluppo:] affronta nel dettaglio la fase realizzativa, argomentando la selezione dei linguaggi di programmazione utilizzati e descrivendo lo sviluppo analitico delle varie componenti del sistema.
    
    \item[\Cref{chap:results} -- Validazione e test:] presenta le prove sperimentali e i test condotti sul campo al fine di verificare la correttezza, la robustezza e le prestazioni dei componenti realizzati. Approfondisce inoltre le strategie e le decisioni operative adottate per il dispiegamento (\emph{deployment}) e l'integrazione delle nuove componenti all'interno dell'infrastruttura di rete e del nuovo applicativo.
    
    \item[\Cref{chap:conclusion} -- Conclusioni:] raccoglie le considerazioni finali sui risultati conseguiti e traccia le possibili direttrici per gli sviluppi e i perfezionamenti futuri.
\end{description}